% H9-T17: exact Zhang scope and negative directions after a finite cover.
% Run JAC1964-20261121; /root/h9_height; GPT-6 Astra Pro.

\subsection*{H9-T17. Threefold heights after a cover: the primitive sign and a negative complement}

\paragraph{The existing characteristic-zero theorem, with both hypotheses.}
Shou-Wu Zhang, \emph{Positivity of heights of codimension 2 cycles
over function field of characteristic 0}, Theorem 1, author pp.~1--2,
states the following. Let $f:\mathcal X\to C$ be flat, with
$\mathcal X,C$ smooth projective and geometrically connected over a
characteristic-zero field, $\dim C=1$, $\dim\mathcal X=n\ge4$,
and let $\mathcal L$ be ample on $\mathcal X$. If a rational
codimension-two cycle $z$ satisfies
\[
 \deg(z\cdot v)=0\quad\text{for every vertical surface }v,
 \qquad
 \operatorname{alb}\bigl(z_Kc_1(\mathcal L_K)^{n-3}\bigr)=0,
                                                               \tag{H9.77}
\]
then
\[
                  \deg(z^2c_1(\mathcal L)^{n-4})\ge0.
                                                               \tag{H9.78}
\]
Equality holds exactly when $z$ is numerically equivalent on
$\mathcal X$ to a fibre of a rational horizontal divisor.
The source was directly read in full at
\url{https://arxiv.org/pdf/1001.4788}; its title page has the
manuscript date 20 September 2021 and the arXiv identifier
\texttt{1001.4788v1}. The positivity theorem is an existing result,
not a deduction claimed as new here.

For a threefold generic fibre $Y/K$, put $h=c_1(\mathcal L_K)$.
The conclusion is $H_Y(z_K,z_K)\ge0$ on the admissible domain
\[
 \mathcal H_h^{\rm Alb}
    =\{z\in\operatorname{CH}^2(Y)_{\rm hom,\mathbf Q}^{\rm adm}:
                  \operatorname{alb}(hz)=0\}.               \tag{H9.79}
\]
The sign is the middle sign $(-1)^2=+1$. The theorem does not
assert that every homologically trivial cycle admits an admissible
extension, and it does not drop the second condition in (H9.77).
On a fixed smooth projective model carrying an ample extension of
$h$, its equality criterion says exactly that the radical on
(H9.79) is
\[
 \mathcal H_h^{\rm Alb}\ \cap\
 \operatorname{im}\bigl(\operatorname{CH}^2(\mathcal X)_{
      \rm num,\mathbf Q}\to\operatorname{CH}^2(Y)_{\mathbf Q}\bigr).
                                                               \tag{H9.80}
\]
Indeed subtract the vertical fibre divisor supplied by Zhang to
obtain a numerically trivial global lift. Conversely such a global
lift is admissible and pairs to zero with every global cycle. A
positive semidefinite rational form has zero vectors exactly in
its radical, so Zhang's equality criterion proves (H9.80).

A height-admissible class may be handled on a smooth projective
model with an ample line bundle whose generic restriction is a
positive multiple of the chosen ample $h$; in characteristic zero
such models are obtained by spreading an embedding, resolving the
vertical locus and twisting relatively ample extensions from the
base. The height and the Albanese-kernel condition are unchanged,
respectively multiplied by positive constants, by that replacement.
Alternatively (H9.77)--(H9.80) can simply be used on a specified
model with its specified ample bundle, as in the example below.

\paragraph{The exact rational Lefschetz splitting of this sign question.}
On any smooth projective threefold $Y/K$, intersection with $h^2$
and Albanese gives an isogeny in the rational category
\[
 \psi_h:\operatorname{Pic}^0(Y)\longrightarrow\operatorname{Alb}_Y,
             \quad D\longmapsto\operatorname{alb}(h^2D).
                                                               \tag{H9.81}
\]
One proof is to take a smooth sufficiently divisible complete-
intersection curve and use the usual Picard/Albanese restriction
and Gysin maps; their composite is an isogeny. This is exactly the
input used in Kahn, \emph{Refined height pairing}, Lemma 6.1,
printed p.~38. After clearing the fixed denominators, the inverse
of (H9.81) acts on rational $K$-points. If necessary a zero-cycle
of positive degree clears the Picard-functor descent obstruction;
in the pencil application there is already a rational point from
the fixed axis.

For an admissible homologically trivial $z$, define
\[
 D_z=\psi_h^{-1}\operatorname{alb}(hz),
       \qquad z_0=z-hD_z.
                                                               \tag{H9.82}
\]
The cycle $hD_z$ is algebraically trivial and is admissible by
Kahn Theorem 5.6. Thus $z_0$ is admissible and belongs to (H9.79).
The decomposition is unique: if $hD$ is in (H9.79), then
$\psi_h(D)=0$, hence $D=0$ in the rational group.

It is an orthogonal decomposition for heights. Projection formula
and Kahn Theorem 6.2(a), which annihilates the Albanese kernel in
the divisor/zero-cycle pairing, give
\[
 H_Y(z_0,hD)=H_Y(hz_0,D)=0,
 \qquad
 H_Y(hD,hD)=H_Y(D,h^2D).                                   \tag{H9.83}
\]
These adjunctions can be expressed using only degree-zero
correspondences: replace $h$ by a positive multiple represented by
a smooth hyperplane surface $i:S\hookrightarrow Y$, and use the
pullback/pushforward height adjunction for $i$ twice. Dividing out
the positive multiples gives (H9.83).
For a curve base the last form is \emph{negative definite} on
\[
 \operatorname{Pic}^0(Y)(K)_{\mathbf Q}/
       \operatorname{im}\operatorname{Tr}_{K/k}
                         (\operatorname{Pic}^0_Y)(k)_{\mathbf Q}
\]
by Kahn Theorem 6.6, printed p.~41. It vanishes on trace points by
Theorem 6.2(c). Consequently in characteristic zero the actual
threefold height has a positive semidefinite Albanese-primitive
part and a negative semidefinite Lefschetz part, with the displayed
exact trace kernel on the latter. It need not be semipositive on
the whole homologically trivial group.

All these statements use admissible cycles. When H9-T12 supplies
admissibility for every homologically trivial cycle, as for a cover
unramified above the critical values of the pencil, this splitting
and the two signs apply to that whole group. They leave the
primitive incidence statement of H9-T12 unchanged.

\paragraph{An explicit negative class on an unramified-above-critical finite cover.}
Work over an algebraically closed field $k$ of characteristic zero.
Let $E/k$ be an elliptic curve and let
\[
 A=E^4,\qquad h_A=d(a_1+a_2+a_3+a_4),
\]
where $a_i$ is the origin divisor from the $i$-th factor and $d$
is large enough to choose a genuine Lefschetz pencil. Let
$f:\widetilde A\to\mathbf P^1$ be its blowup model and let
$\Sigma$ be its finite set of critical values.

The quotient map $E\to E/\{\pm1\}\simeq\mathbf P^1$ is separable
of degree two and has four branch values. Postcompose it with a
projective linear transformation so that all four branch values
avoid $\Sigma$; such a transformation exists because $k$ is
infinite and the forbidden conditions form a finite union of
proper closed subsets of $\operatorname{PGL}_2$. Call the result
\[
                   \pi:C=E\longrightarrow\mathbf P^1.
\]
It is unramified above every critical value. Put
$K=k(\mathbf P^1)$, $L=k(C)$, and
$\mathcal X=\widetilde A\times_{\mathbf P^1}C$.
This total space is smooth: above $\Sigma$ the base change is
etale, while above a branch value the original morphism is smooth.
It is projective and flat over $C$, and its generic fibre $Y_L$
is the base change of the pencil threefold.

The pencil incidence realizes $\mathcal X$ as a divisor in
$A\times C$ of class
\[
                        h_A+b,\qquad b=\pi^*c_1(\mathcal O(1)).
\]
The product ample bundle with class $h_A+b$ restricts to an ample
bundle $\mathcal L$ on $\mathcal X$ whose generic restriction is
$h=h_A|_{Y_L}$. Thus the example uses exactly the ample-model
setting of Zhang's theorem.

On $C\times E_1$, where $E_1$ is the first factor of $A$, take
\[
 \mathcal P=\mathcal O(\Delta-C\times\{0\}-\{0\}\times E_1).
\]
Pull it to $A\times C$ and then to $\mathcal X$; denote the
resulting divisor class by $\mathcal D$, and its generic
restriction by $D$. The associated point of
$\operatorname{Pic}^0(A)(L)$ is the nonconstant identity map from
$C=E$ into its first dual elliptic factor. Restriction identifies
$\operatorname{Pic}^0(A)$ with $\operatorname{Pic}^0(Y_L)$ by the
Lefschetz theorem (or its integral degree-one cohomology comparison
and the Picard complex torus description). In particular $D$ is
nonzero modulo the constant $L/k$-trace.

Put
\[
                z=hD\in\operatorname{CH}^2(Y_L)_{\mathbf Q},
       \qquad \widetilde z=c_1(\mathcal L)\mathcal D.
\]
Every fibre of $\mathcal D$ is the restriction of an algebraically
trivial line bundle on $A$, hence has zero rational cycle class on
all vertical cycles. Therefore $\widetilde z$ is an admissible
extension directly. The class $z$ is algebraically, and hence
geometrically homologically, trivial.

Nevertheless its height is strictly negative. One may either use
(H9.83) and Kahn's negativity theorem, or compute it exactly on the
smooth total space:
\[
 \begin{aligned}
 H_{Y_L}(z,z)
  &=\int_{\mathcal X}c_1(\mathcal L)^2\mathcal D^2\\
  &=\int_{A\times C}(h_A+b)^3\mathcal P^2\\
  &=\int_{A\times C}h_A^3\mathcal P^2
    =6d^3\int_{C\times E_1}\mathcal P^2
    =-12d^3<0.
 \end{aligned}                                               \tag{H9.84}
\]
Here $\mathcal P^2$ already has top codimension on $C\times E_1$,
so every term containing $b$ or $a_1$ vanishes. Only
$6d^3a_2a_3a_4$ remains from $h_A^3$. Finally
$\deg_{C\times E_1}\mathcal P^2=-2$, since the diagonal and the
two axes each have self-intersection zero, and each distinct pair
has intersection one. No conjectural positivity statement or
numerical approximation is used in (H9.84).

This negative class already lies in the new norm-zero summand.
The deck involution of $C/\mathbf P^1$ is $[-1]$.
The normalized Poincare bundle satisfies
$([-1]\times1)^*\mathcal P=\mathcal P^{-1}$, while $h$ is
invariant. Hence $\iota^*z=-z$. The Galois push-pull formula gives
\[
 \operatorname{res}\operatorname{cor}(z)=z+\iota^*z=0.
\]
Rational restriction is injective because
$\operatorname{cor}\operatorname{res}=2$, so
\[
                         \operatorname{cor}(z)=0.            \tag{H9.85}
\]
No further subtraction of the old restriction component is needed.

\paragraph{The negative example in every characteristic.}
The construction and the exact number (H9.84) also work over an
algebraically closed field of arbitrary characteristic. In fact the
same involution construction works: $[-1]$ is nontrivial on every
elliptic curve, since $[2]$ is not the zero morphism, and the fixed
field of this order-two automorphism gives a separable Galois
degree-two extension even in characteristic two. The quotient is
$\mathbf P^1$ by the Weierstrass $x$ coordinate. Its finite branch
set may be wild and need not have four points; moving that set
off $\Sigma$ is all the smoothness argument needs. Rational
push-pull and the anti-invariance proof (H9.85) remain valid.
In positive characteristic Picard restriction need only be used as
an isogeny, which suffices after tensoring with $\mathbf Q$.

There is also a construction avoiding degree divisible by the
characteristic. To obtain
a convenient separable cover, choose an integer $N\ge3$ prime to
the characteristic and a basepoint-free pencil in
$|\mathcal O_E(N\cdot0)|$. Such a pencil exists because this line
bundle is very ample and two general sections have disjoint zero
divisors. It gives a degree-$N$ morphism $E\to\mathbf P^1$.
The inseparable degree divides both a power of the characteristic
and $N$, so the map is separable. Postcompose by a projective
linear transformation to move its finite branch set away from
$\Sigma$. The same smooth incidence model and admissible class
$hD$ then give $H(hD,hD)=-12d^3$.

There is still exact norm zero, even though this cover need not be
Galois. Every scheme-theoretic fibre divisor of $\pi$ is linearly
equivalent to $N\cdot0$. Its sum in the elliptic group, including
multiplicities, is therefore zero. Norm of the normalized Poincare
line bundle sends the generic parameter point to precisely this
sum. Hence $\operatorname{cor}(D)=0$ in
$\operatorname{Pic}^0(A_K)$; norm commutes with restriction to
$Y_K$, and projection formula gives
$\operatorname{cor}(hD)=h\operatorname{cor}(D)=0$.
No assertion of tame ramification is needed: all branch values
were placed in the smooth locus of the pencil. This proves the
negative norm-zero example in every characteristic in which the
chosen genuine Lefschetz pencil is taken; sufficiently ample
pencils on the abelian fourfold exist in all characteristics.
Zhang's positive Albanese-primitive conclusion above remains a
characteristic-zero assertion unless its additional Hodge-index
input is supplied.

\paragraph{An explicit positive direction in the same new norm-zero group.}
In the same example put $\alpha=a_2-a_3$ on $A$ and
$z_+=\alpha|_{Y_L}D$, while $z_-=hD$ is the previous negative
class. Both are algebraically trivial and admissible, and both
have corestriction zero by the same Poincare norm calculation.
On the incidence model, the intersection identities are
\[
 H(z_+,z_+)=4d>0,\qquad H(z_-,z_-)=-12d^3<0,
                   \qquad H(z_+,z_-)=0.                    \tag{H9.85a}
\]
Indeed $\alpha^2=-2a_2a_3$, and only $d a_4$ from $h_A$
survives after multiplication by $\mathcal P^2$; this gives
$(-2d)(-2)=4d$. For the mixed pairing the coefficient of
$a_2a_3a_4$ in $h_A^2(a_2-a_3)$ is zero, and all other terms
vanish against $\mathcal P^2$.

This positive direction satisfies the Albanese-primitive condition.
On $A$, the exact Chow identity
$h_A^2\alpha D_A=0$ holds: terms containing $a_1D_A$ vanish on
the first elliptic factor, and the remaining two triple terms
cancel. Thus $j_*(h_Yz_+)=0$; the Albanese map induced by the
hyperplane inclusion $j:Y_L\hookrightarrow A_L$ is an isogeny,
so $\operatorname{alb}(h_Yz_+)=0$ rationally. Its family
realization is in the primitive middle factor as well, since
$h_A^2(\alpha\,H^1(E_1))=0$.
Consequently (H9.85a) exhibits a rational plane of signature
$(1,1)$ inside the new kernel of norm in every characteristic.
It is a concrete positive primitive direction alongside the
negative Lefschetz direction, not a proof of positivity on all
primitive classes.

\paragraph{The three notions of primitivity must be kept distinct.}
The cycle in (H9.84) has $\operatorname{cl}(z)=0$, so every
cohomological primitive equation applied to its ordinary cycle
class holds vacuously. It is not Chow primitive: $hz=h^2D$ has
nonzero Albanese class, because (H9.81) is an isogeny and $D$ is
nonconstant modulo trace. In particular $hz\ne0$ in rational Chow.
It also fails precisely Zhang's condition
$\operatorname{alb}(hz)=0$. On the incidence realization it belongs
to the Lefschetz part, not the primitive middle factor of H9-T9.

Thus this example refutes semipositivity on the \emph{entire}
homologically trivial group after an arbitrary separable cover,
even in characteristic zero and even when the cover is unramified
above all critical values. It does not contradict H9-T7 over the
original rational pencil field: there the Picard variety is
constant and $\operatorname{Hom}(J(\mathbf P^1),\operatorname{Pic}^0A)=0$,
so this nonconstant Lefschetz direction does not exist. It does not
contradict the finite-extension primitive polarization/trace theorem
of H9-T12. The corrected finite-cover statement is the orthogonal
positive-primitive/negative-Lefschetz decomposition above, on its
explicit admissible domain.

\paragraph{Sources and check scope.}
Zhang's complete four-page primary note and Kahn's Lemma 6.1,
Theorems 6.2 and 6.6 were directly read. The new contributions here
are the exact scope comparison, rational decomposition (H9.82),
and the explicit unramified-above-critical counterexample with
intersection (H9.84) and norm identity (H9.85). General height
functoriality is Kahn Theorem 4.14. No historical originality or
formal proof-assistant verification is claimed.
